
\begin{abstract}
%\shafiq{Let's put a project page}


%\zixuan{Title: MAS-Orchestra: Elevating Collective Intelligence via Multi-Agent System Orchestration — A Training Framework, Benchmark, and Analysis?}

% \zixuan{working on Abstract}

While multi-agent systems (MAS) promise elevated intelligence through coordination of agents, current approaches to automatic MAS design under-deliver. Such shortcomings stem from two key factors: (1) methodological complexity -- agent orchestration is performed using sequential, code-level execution that limits global system-level holistic reasoning and scales poorly with agent complexity -- and (2) efficacy uncertainty -- MAS are deployed without understanding if there are tangible benefits compared to single-agent systems (SAS). We propose \ourframework, a training-time framework that formulates MAS orchestration as a function-calling reinforcement learning problem with holistic orchestration, generating an entire MAS at once. In \ourframework, complex, goal-oriented sub-agents are abstracted as callable functions, enabling global reasoning over system structure while hiding internal execution details. To rigorously study when and why MAS are beneficial, we introduce \ourbenchmark, a controlled benchmark that characterizes tasks along five axes: \depth, \horizon, \breadth, \myparallel, and \robustness. Our analysis reveals that MAS gains depend critically on task structure, verification protocols, and the capabilities of both orchestrator and sub-agents, rather than holding universally. Guided by these insights, \ourframework achieves consistent improvements on public benchmarks including mathematical reasoning,  multi-hop QA, and search-based QA, {while achieving more than \textit{10$\times$} efficiency over strong baselines}. Together, \ourframework and \ourbenchmark enable better training and understanding of MAS in the pursuit of multi-agent intelligence.
% \footnote{We will open-source the data, checkpoint, code, leaderboard for all components upon acceptance.}
% \footnote{Code, data, examples and the trained orchestrators will be publicly available upon the acceptance}




%Multi-agent systems (MAS) promise stronger collective intelligence by enabling coordination among agents. However, existing approaches to automatic MAS design often formulate orchestration as executable code that does not scale to complex agents, rely on sequential decisions that prevent global system-level reasoning, and lack principled understanding of when multi-agent coordination is beneficial over single-agent systems (SAS).



%We propose \ourframework, a training-time framework that formulates MAS orchestration as a function-calling reinforcement learning problem with holistic orchestration, generating a complete MAS in a single step. This formulation encapsulates complex, goal-oriented sub-agents as callable functions, enabling global reasoning over system structure without exposing sub-agent internals. To systematically analyze MAS benefits over SAS, we introduce \ourbenchmark, a controlled benchmark that characterizes tasks along five axes: \depth, \horizon, \breadth, \myparallel, and \robustness. Using this benchmark, we conduct rigorous analysis that reveals MAS effectiveness depends on task structure, verification protocols, and the capabilities of both the orchestrator and sub-agents, rather than holding universally. Guided by these insights, \ourframework achieves consistent improvements on public benchmarks, including mathematical reasoning, multi-hop question answering, and search-based question answering. 
%Together, \ourframework, \ourbenchmark, and comprehensive analysis provide a principled foundation for designing and understanding MAS that genuinely elevate collective intelligence.


%\austin{Abstract is a bit lengthy, 1/18 attempt to condense}

 %This work addresses both factors: We propose \ourframework, a simple and flexible training framework that formulates agent orchestration as a function-calling reinforcement learning (RL) problem. The formulation of goal-oriented sub-agents as callable functions enables global reasoning over system structure.  We also propose \ourbenchmark, a benchmark engine that enables analysis of MAS and SAS over five axes: \depth, \horizon, \breadth, \myparallel, \austin{Parallel maybe should be Parallelism (Noun)} and \robustness. Using \ourbenchmark, we rigorously analyze specific factors such as task structure, verification protocols, and orchestrator/sub-agent interplay to guide our MAS development. In all, MAS trained with \ourframework yield consistent improvements in math reasoning, multi-hop question answering, and search-based question answering. 
 
 
\end{abstract}
